\textbf{Elección y deducción del procedimiento.} El perfil responde al ángulo efectivo: $C_L=a_0[\alpha-\alpha_0-C_L/(\pi e_i AR)]$. Agrupar términos en $C_L$ da $C_L[1+a_0/(\pi e_i AR)]=a_0(\alpha-\alpha_0)$, de donde se lee la pendiente del ala.

\begin{enumerate}

\item La corrección por ángulo inducido aumenta el denominador.

\[\fitmath{\mathrm{factor}=1+\frac{a0}{\pi\cdot \mathrm{ei}\cdot \mathrm{AR}}}\]

\[\fitmath{\mathrm{factor}=1+\frac{(5.9)}{\pi\cdot (0.91)\cdot (8.2)}\approx 1.251679\,}\]

\item La pendiente se interpreta por radián.

\[\fitmath{a=\frac{a0}{\mathrm{factor}}}\]

\[\fitmath{a=\frac{(5.9)}{(1.25167895)}\approx 4.713669\,}\]

\end{enumerate}
